\documentclass[11pt,a4paper]{article}
\usepackage{fontspec,amsmath,amssymb,amsthm,microtype}
\setmainfont{lmroman10-regular.otf}[BoldFont=lmroman10-bold.otf,ItalicFont=lmroman10-italic.otf,BoldItalicFont=lmroman10-bolditalic.otf]
\usepackage[margin=25mm]{geometry}
\usepackage[unicode,hidelinks]{hyperref}
\hypersetup{pdfauthor={LangevinLean Project}}
\setlength{\emergencystretch}{3em}
\setlength{\parskip}{0.4em}
\setlength{\parindent}{0pt}
\allowdisplaybreaks[2]
\raggedbottom
\title{\texorpdfstring{Exact Hamiltonian energy and volume-\allowbreak{}preserving Verlet geometry}{Exact Hamiltonian energy and volume-preserving Verlet geometry}}
\author{LangevinLean Project}
\date{}
\begin{document}\maketitle
Draft mathematical exposition; these notes may contain mistakes.


The Hamiltonian subflow is the kinetic equation with zero friction and zero noise. Its total energy is conserved by an exact cancellation between force work and kinetic-energy change. This section proves conservation for the Picard-constructed integral flow and proves that the explicit kick/\allowbreak{}drift/\allowbreak{}Verlet maps preserve product Haar volume. These are prerequisites for Gibbs stationarity; energy conservation alone and volume preservation of a numerical map do not yet establish invariance of the full kinetic diffusion.



Exact energy conservation belongs to the continuous Hamiltonian flow. Verlet instead supplies reversible volume-preserving geometry with a nonzero energy defect, making it suitable for a corrected invariant comparison.


\section{\texorpdfstring{Differentiate total energy along the Hamiltonian equations}{Differentiate total energy along the Hamiltonian equations}}


Define \(H(v,x)=f(x)+\|v\|^2/2\). If \(x'=v\) and \(v'=-\nabla f(x)\), the chain rule gives \((f(x))'=\langle\nabla f(x),v\rangle\), while the squared-norm derivative gives \((\|v\|^2/2)'=-\langle v,\nabla f(x)\rangle\). Symmetry of the real inner product makes their sum zero. The proof composes the Fréchet derivative of \(f\) with the position derivative, uses the gradient representation of that derivative, and combines it with the Hilbert squared-norm derivative. The statement uses right derivatives within \([t,\infty)\), allowing the initial endpoint of an integral solution.



For \(H(v,x)=\|v\|^2/2+f(x)\), the two derivative contributions are \(\langle v,-\nabla f(x)\rangle\) and \(\langle\nabla f(x),v\rangle\). They cancel exactly. This calculation concerns the true Hamiltonian differential equation, before any numerical discretization is introduced.


\section{\texorpdfstring{Turn vanishing right derivatives into exact conservation}{Turn vanishing right derivatives into exact conservation}}


A continuous scalar function with zero right derivative on \([0,T)\) is constant on \([0,T]\), including the final endpoint. The proof applies the scalar Gronwall comparison with coefficient and forcing both zero, first to \(F\) and then to \(-F\). These give \(F(t)\le F(0)\) and \(-F(t)\le-F(0)\). Their combination proves equality. This avoids assuming a derivative of the chosen extension beyond the terminal time or on negative time.



A continuous scalar function with zero right derivative on the interval is constant. Applying this to the composed Hamiltonian proves energy conservation up to the endpoints. The argument does not require differentiability of the flow at negative times.


\section{\texorpdfstring{Conserve energy for the constructed Hamiltonian integral flow}{Conserve energy for the constructed Hamiltonian integral flow}}


the Hamiltonian flow is the already constructed the kinetic solution map with friction zero, inverse mass one, and the zero compact forcing path. Its specification reduces to \(Y_t=Y_0-\int_0^tD(Y_s)ds\), with \(D(v,x)=(\nabla f(x),-v)\). Continuity of the Lipschitz drift gives the integral derivative. The argument identifies the selected path with its differentiable primitive on a right neighborhood, projects its derivative onto velocity and position, and applies the energy derivative cancellation. The scalar constancy theorem then yields \(H(Y_t)=H(Y_0)\). This is an exact statement about an existing solution, not an energy-preservation hypothesis on an abstract flow.



The selected Hamiltonian solution satisfies the integral equations with zero friction and zero forcing. Differentiating those equations supplies the hypotheses of the conservation calculation. Consequently conservation belongs to the constructed flow, rather than to an unspecified trajectory satisfying a formal equation.


\section{\texorpdfstring{Build the reversible proposal and preserve the Hamiltonian under reversal}{Build the reversible proposal and preserve the Hamiltonian under reversal}}


Velocity reversal \(R(v,x)=(-v,x)\) conjugates every kick and drift to its negative-time inverse. Applying this to the palindromic kick–drift–kick composition gives \(R\,V_h\,R=V_{-h}\). Since \(V_{-h}V_h=\mathrm{id}\), the proposal \(J_h=R\,V_h\) satisfies \(J_h^2=\mathrm{id}\). The argument rewrites the three component conjugacies and then uses the proved inverse Verlet identity. Negating velocity leaves its squared norm unchanged, so \(H(Rz)=H(z)\). The exact involution will give symmetric old/\allowbreak{}new proposal graph measures in the Metropolis argument.



Velocity reversal preserves the kinetic square and leaves position fixed. Composing it with a reversible Verlet map creates an involutive proposal. An involution is useful for Metropolis balance because its forward and backward proposal graphs are paired by the same transformation.


\section{\texorpdfstring{Expose the exact numerical energy defect}{Expose the exact numerical energy defect}}


Let \(w=v-hg(x)/2\), \(y=x+hw\), and \(v^+=w-hg(y)/2\). Expanding both kinetic-energy squares gives

\[H(v^+,y)-H(v,x)=f(y)-f(x)-\frac h2\langle w,g(x)+g(y)\rangle
+\frac{h^2}{8}\bigl(\|g(y)\|^2-\|g(x)\|^2\bigr).\]

The first part is a trapezoidal force-work remainder. The second is an explicit force-norm correction. First, the argument proves the two-half-kick norm identity for arbitrary vectors using squared-norm and inner-product expansions, then substitutes the intermediate Verlet velocity and position. This identity preserves the error terms that must later be bounded to control rejection; no second-order error estimate is assumed by definition.



The numerical energy change is obtained by evaluating the true Hamiltonian before and after one Verlet step. It generally does not vanish. Keeping it explicit permits both a Metropolis acceptance formula and quantitative rejection estimates derived from force smoothness.


\section{\texorpdfstring{Prove the elementary subflows preserve Haar volume}{Prove the elementary subflows preserve Haar volume}}


A drift \((v,x)\mapsto(v,x+hv)\) is a skew product whose second component is a translation for each fixed velocity. Translation invariance of Haar measure therefore proves preservation of the product measure. A kick \((v,x)\mapsto(v-hg(x),x)\) is the same construction after swapping the two phase coordinates; measurability of \(g\) suffices. The argument applies the measure-preserving skew-product theorem to translations, then conjugates the kick by the measure-preserving swap. Composition proves the complete Verlet map preserves product Haar volume. This derives the Jacobian-one property without presupposing differentiability of the force.



Each kick is a translation of velocity depending only on position; each drift is a translation of position depending only on velocity. Fubini and translation invariance show that these shears preserve product volume. Their compositions and velocity reversal preserve volume as well.


\section{\texorpdfstring{The reversed Verlet proposal is a measure-\allowbreak{}preserving involution}{The reversed Verlet proposal is a measure-preserving involution}}


For a Haar factor also invariant under negation, velocity reversal preserves product measure: apply the negation-preserving map in the velocity coordinate and identity in the position coordinate. Composing this with the proved volume-preserving Verlet map gives volume preservation of the explicit proposal \(J_h\). Together with \(J_h^2=\mathrm{id}\), this supplies the two exact properties used by deterministic Metropolis detailed balance. Neither property asserts that the uncorrected finite-step Verlet map itself preserves the Gibbs weight.



The reversed Verlet proposal combines volume preservation with involutivity. These are exactly the geometric ingredients needed for the accepted-flow symmetry argument. Approximate energy conservation is not needed to prove corrected invariance, although it controls how often proposals are rejected.


\end{document}
